% ============================================================
%  Chapter 3 -- Series, Nonnegative Terms, e: 3.21-3.32
% ============================================================
\rsection{Series}

\begin{signpost}[A series is a sequence in disguise]
Nothing new is defined in this section except notation: a series \emph{is} its
sequence of partial sums, and Rudin says so twice --- in 3.21 (``$s$ is the
limit of a sequence of sums, and is not obtained simply by addition'') and in
the remark that every theorem about sequences can be restated for series via
$a_n = s_n - s_{n-1}$. The three theorems here are the three big sequence facts
passed through that dictionary: 3.22 is the Cauchy criterion 3.11, 3.24 is
monotone convergence 3.14, and 3.23 is the trivial-but-vital corollary that
terms of a convergent series must die out.

\smallskip
The section ends with the comparison test, 3.25 --- the master test. Every
test that follows (root, ratio, condensation applications) works by comparing
against a series already understood, almost always the geometric series 3.26.
\end{signpost}

\noindent
In the remainder of this chapter, all sequences and series under consideration
will be complex-valued, unless the contrary is explicitly stated. Extensions of
some of the theorems which follow, to series with terms in $\R^k$, are
mentioned in Exercise 3.15.

\ritem{3.21}{Definition}
Given a sequence $(a_n)$, we use the notation
\[ \sum_{n=p}^{q} a_n \qquad (p \le q) \]
to denote the sum $a_p + a_{p+1} + \cdots + a_q$. With $(a_n)$, we associate a
sequence $(s_n)$, where
\[ s_n = \sum_{k=1}^{n} a_k. \]
For $(s_n)$ we also use the symbolic expression
\[ a_1 + a_2 + a_3 + \cdots \]
or, more concisely,
\begin{equation}
  \sum_{n=1}^{\infty} a_n. \tag{3.4}
\end{equation}
The symbol in (3.4) we call an \emph{infinite series}, or just a \emph{series}.
The numbers $s_n$ are called the \emph{partial sums} of the series. If $(s_n)$
converges to $s$, we say that the series \emph{converges}, and write
\[ \sum_{n=1}^{\infty} a_n = s. \]
The number $s$ is called the \emph{sum} of the series; but it should be clearly
understood that $s$ is the limit of a sequence of sums, and is not obtained
simply by addition.\kw{This sentence is the guard rail for the whole second
half of the chapter. Addition of finitely many numbers is commutative and
associative by the field axioms; a ``sum'' that is really a \emph{limit}
inherits neither property automatically. Theorem 3.54 will show rearrangement
genuinely failing, and 3.55 will show what hypothesis buys it back.}

If $(s_n)$ diverges, the series is said to \emph{diverge}.

Sometimes, for convenience of notation, we shall consider series of the form
\begin{equation}
  \sum_{n=0}^{\infty} a_n. \tag{3.5}
\end{equation}
And frequently, when there is no possible ambiguity, or when the distinction is
immaterial, we shall simply write $\Sigma a_n$ in place of (3.4) or (3.5).

It is clear that every theorem about sequences can be stated in terms of series
(putting $a_1 = s_1$ and $a_n = s_n - s_{n-1}$ for $n > 1$), and vice versa.
But it is nevertheless useful to consider both concepts.

The Cauchy criterion (Theorem 3.11) can be restated in the following form:

\ritem{3.22}{Theorem}
\emph{$\Sigma a_n$ converges if and only if for every $\varepsilon > 0$ there
is an integer $N$ such that}
\begin{equation}
  \left|\sum_{k=n}^{m} a_k\right| \le \varepsilon \tag{3.6}
\end{equation}
\emph{if $m \ge n \ge N$.}

In particular, by taking $m = n$, (3.6) becomes
\[ |a_n| \le \varepsilon \qquad (n \ge N). \]
In other words:

\ritem{3.23}{Theorem}
\emph{If $\Sigma a_n$ converges, then $\lim_{n\to\infty} a_n = 0$.}

The condition $a_n \to 0$ is not, however, sufficient to ensure convergence of
$\Sigma a_n$. For instance, the series
\[ \sum_{n=1}^{\infty} \frac{1}{n} \]
diverges; for the proof, we refer to Theorem 3.28.\kw{Necessary, not
sufficient --- the single most-cited warning in elementary analysis. The
harmonic series is the canonical witness: terms shrink to $0$, partial sums
climb past every bound. Rudin flags it before anyone can form the converse
habit; 3.22 with $m > n$ is what genuine convergence requires, and a
\emph{block} of small terms can still add up.}

Theorem 3.14, concerning monotonic sequences, also has an immediate
counterpart for series.

\ritem{3.24}{Theorem}
\emph{A series of nonnegative\footnote{The expression ``nonnegative'' always
refers to real numbers.} terms converges if and only if its partial sums form a
bounded sequence.}\kn{Nonnegative terms make the partial sums monotonically
increasing, and 3.14 does the rest. This innocent-looking equivalence is the
engine of the whole nonnegative theory: comparison (3.25), condensation
(3.27), and the root test all reduce convergence to boundedness of partial
sums, which is checked by comparing.}

We now turn to a convergence test of a different nature, the so-called
``comparison test.''

\ritem{3.25}{Theorem}
\begin{enumerate}[label=(\alph*), leftmargin=2.2em, itemsep=2pt, topsep=3pt]
  \item \emph{If $|a_n| \le c_n$ for $n \ge N_0$, where $N_0$ is some fixed
        integer, and if $\Sigma c_n$ converges, then $\Sigma a_n$ converges.}
  \item \emph{If $a_n \ge d_n \ge 0$ for $n \ge N_0$ and if $\Sigma d_n$
        diverges, then $\Sigma a_n$ diverges.}
\end{enumerate}
\emph{Note that (b) applies only to series of nonnegative terms $a_n$.}

\begin{proof}
Given $\varepsilon > 0$, there exists $N \ge N_0$ such that $m \ge n \ge N$
implies
\[ \sum_{k=n}^{m} c_k \le \varepsilon, \]
by the Cauchy criterion. Hence
\[ \left|\sum_{k=n}^{m} a_k\right| \le \sum_{k=n}^{m} |a_k|
   \le \sum_{k=n}^{m} c_k \le \varepsilon, \]
and (a) follows.

Next, (b) follows from (a), for if $\Sigma a_n$ converges, so must $\Sigma d_n$
[note that (b) also follows from Theorem 3.24].
\end{proof}

\begin{unpack}[What comparison actually proves, and the asymmetry]
Look at the chain of inequalities: it bounds $|\sum a_k|$ through
$\sum|a_k|$. So (a) silently proves more than it claims --- under its
hypotheses $\Sigma|a_n|$ converges too, i.e.\ the convergence is
\emph{absolute} (3.45's terminology). This is why Rudin will say at 3.46 that
comparison, root and ratio ``cannot give any information about conditionally
convergent series'': all three run through this chain.

The asymmetry between (a) and (b) is not decoration. Domination transfers
convergence \emph{downward} and divergence \emph{upward}, never the reverse;
and (b) needs $a_n \ge 0$ because a wildly oscillating series can dominate a
divergent positive one and still converge --- cancellations rescue it. Keep
the mantra: below a convergent series, convergent; above a divergent
nonnegative series, divergent.
\end{unpack}

The comparison test is a very useful one; to use it efficiently, we have to
become familiar with a number of series of nonnegative terms whose convergence
or divergence is known.

\rsection{Series of Nonnegative Terms}

\noindent
The simplest of all is perhaps the geometric series.

\ritem{3.26}{Theorem}
\emph{If $0 \le x < 1$, then}
\[ \sum_{n=0}^{\infty} x^n = \frac{1}{1-x}. \]
\emph{If $x \ge 1$, the series diverges.}

\begin{proof}
If $x \ne 1$,
\[ s_n = \sum_{k=0}^{n} x^k = \frac{1-x^{n+1}}{1-x}. \]
The result follows if we let $n \to \infty$. For $x = 1$, we get
\[ 1 + 1 + 1 + \cdots, \]
which evidently diverges.\kn{One of only two series in the chapter whose sum
is \emph{computed} rather than merely shown to exist (the other is $e$). The
closed form for $s_n$ --- multiply by $x$, subtract, telescope --- is the
entire proof, and the limit uses 3.20(e). Every later test converges things by
comparing them to this one series; it is the yardstick.}
\end{proof}

\begin{center}
\begin{tikzpicture}[x=1mm, y=1mm]
  \draw[axis] (-2,0) -- (95,0);
  \foreach \x in {0,36,57.6,70.56,78.34,83}{\draw[thin] (\x,-1.6) -- (\x,1.6);}
  \node[ptopen] at (90,0) {};
  \node[mlbl, anchor=north] at (18,-2.2)  {$1$};
  \node[mlbl, anchor=north] at (46.8,-2.2){$x$};
  \node[mlbl, anchor=north] at (64.1,-2.2){$x^2$};
  \node[mlbl, anchor=north] at (74.5,-2.2){$x^3$};
  \node[mlbl, anchor=north] at (85.8,-2.2){$\cdots$};
  \draw[thin, {Stealth[length=2.5pt]}-{Stealth[length=2.5pt]}]
    (0,6.5) -- (89.4,6.5);
  \node[mlbl, anchor=south] at (44.7,7.5) {$S$};
  \draw[guide, {Stealth[length=2.5pt]}-{Stealth[length=2.5pt]}]
    (36.4,-7) -- (89.4,-7);
  \node[mlbl, anchor=north] at (62.9,-8) {$xS$ --- the whole, scaled by $x$};
\end{tikzpicture}
\end{center}
\figcap{The segments $1, x, x^2, \dots$ laid end to end (drawn for $x = 0.6$).
The tail beyond the first segment is a copy of the entire picture shrunk by
the factor $x$, so $S = 1 + xS$, i.e.\ $S = 1/(1-x)$. Rudin's closed form for
$s_n$ is this picture stopped after $n$ steps --- it is the identity of 1.21,
$b^{n+1}-a^{n+1} = (b-a)(b^n + \cdots + a^n)$, with $b = 1$, $a = x$.}

In many cases which occur in applications, the terms of the series decrease
monotonically. The following theorem of Cauchy is therefore of particular
interest. The striking feature of the theorem is that a rather ``thin''
subsequence of $(a_n)$ determines the convergence or divergence of
$\Sigma a_n$.

\ritem{3.27}{Theorem}
\emph{Suppose $a_1 \ge a_2 \ge a_3 \ge \cdots \ge 0$. Then the series
$\sum_{n=1}^{\infty} a_n$ converges if and only if the series}
\begin{equation}
  \sum_{k=0}^{\infty} 2^ka_{2^k} = a_1 + 2a_2 + 4a_4 + 8a_8 + \cdots
  \tag{3.7}
\end{equation}
\emph{converges.}

\begin{proof}
By Theorem 3.24, it suffices to consider the boundedness of the partial sums.
Let
\begin{align*}
  s_n &= a_1 + a_2 + \cdots + a_n,\\
  t_k &= a_1 + 2a_2 + \cdots + 2^ka_{2^k}.
\end{align*}
For $n < 2^k$,
\begin{align*}
  s_n &\le a_1 + (a_2+a_3) + \cdots + (a_{2^k}+\cdots+a_{2^{k+1}-1})\\
      &\le a_1 + 2a_2 + \cdots + 2^ka_{2^k}\\
      &= t_k,
\end{align*}
so that
\begin{equation}
  s_n \le t_k. \tag{3.8}
\end{equation}
On the other hand, if $n > 2^k$,
\begin{align*}
  s_n &\ge a_1 + a_2 + (a_3+a_4) + \cdots +
        (a_{2^{k-1}+1}+\cdots+a_{2^k})\\
      &\ge \tfrac12a_1 + a_2 + 2a_4 + \cdots + 2^{k-1}a_{2^k}\\
      &= \tfrac12t_k,
\end{align*}
so that
\begin{equation}
  2s_n \ge t_k. \tag{3.9}
\end{equation}
By (3.8) and (3.9), the sequences $(s_n)$ and $(t_k)$ are either both bounded
or both unbounded. This completes the proof.
\end{proof}

\begin{center}
\begin{tikzpicture}[x=1mm, y=1mm]
  % bars for a_2..a_15 of a decreasing series, grouped in dyadic blocks
  % heights ~ 12/n
  \foreach \n/\h in {2/6.0, 3/4.0}{
    \draw[thin, fill=black!14] ({(\n-2)*5.6},0) rectangle ({(\n-1)*5.6-0.7},\h);}
  \foreach \n/\h in {4/3.0, 5/2.4, 6/2.0, 7/1.71}{
    \draw[thin, fill=black!26] ({(\n-2)*5.6},0) rectangle ({(\n-1)*5.6-0.7},\h);}
  \foreach \n/\h in {8/1.5, 9/1.33, 10/1.2, 11/1.09, 12/1.0, 13/0.92, 14/0.86, 15/0.8}{
    \draw[thin, fill=black!8] ({(\n-2)*5.6},0) rectangle ({(\n-1)*5.6-0.7},\h);}
  % group braces (drawn as brackets below)
  \draw[thin] (0,-2) -- (0,-3.5) -- (10.5,-3.5) -- (10.5,-2);
  \node[lbl, anchor=north] at (4,-4) {$a_2{+}a_3 \le 2a_2$};
  \draw[thin] (11.2,-2) -- (11.2,-8.5) -- (33,-8.5) -- (33,-2);
  \node[lbl, anchor=north] at (23,-9) {$a_4{+}\cdots{+}a_7 \le 4a_4$};
  \draw[thin] (33.6,-2) -- (33.6,-3.5) -- (77.7,-3.5) -- (77.7,-2);
  \node[lbl, anchor=north] at (55.6,-4) {$a_8{+}\cdots{+}a_{15} \le 8a_8$};
  \node[lbl, anchor=west] at (80,8) {each block trapped};
  \node[lbl, anchor=west] at (80,4.5) {between $\tfrac12\,2^ka_{2^k}$};
  \node[lbl, anchor=west] at (80,1) {and $2^ka_{2^k}$};
\end{tikzpicture}
\end{center}
\figcap{Cauchy condensation: chop the series into blocks of doubling length.
Monotonicity makes each block at most its first term times its length, and at
least its last term times half its length --- inequalities (3.8) and (3.9).}

\begin{deeper}[Why the sampling is allowed to be so thin]
The theorem replaces all of $(a_n)$ by the sample $a_1, a_2, a_4, a_8, \dots$
--- vanishingly sparse --- and loses nothing. Monotonicity is what pays for
it: within a dyadic block the terms cannot vary by more than the endpoints, so
the whole block is measured by one sample to within a factor of $2$. Without
monotonicity the sampled series says nothing (put all the mass off the sample
points).

The choice of base $2$ is a convenience; any factor works. And the theorem is
built for exactly one purpose, visible in 3.28--3.29: it turns
$\sum 1/n^p$ into a \emph{geometric} series, and $\sum 1/[n(\log n)^p]$ into
$\sum 1/k^p$ --- each application of condensation strips off one logarithm.
That is why the borderline series in Rudin's ladder all involve iterated logs.
\end{deeper}

\ritem{3.28}{Theorem}
\emph{$\sum \dfrac{1}{n^p}$ converges if $p > 1$ and diverges if $p \le 1$.}

\begin{proof}
If $p \le 0$, divergence follows from Theorem 3.23. If $p > 0$, Theorem 3.27
is applicable, and we are led to the series
\[ \sum_{k=0}^{\infty} 2^k \cdot \frac{1}{2^{kp}} =
   \sum_{k=0}^{\infty} 2^{(1-p)k}. \]
Now, $2^{1-p} < 1$ if and only if $1 - p < 0$, and the result follows by
comparison with the geometric series (take $x = 2^{1-p}$ in Theorem
3.26).\kn{With $p = 1$ this settles the harmonic series: condensation turns it
into $\sum_k 2^k\cdot 2^{-k} = 1+1+1+\cdots$. The classical grouped proof ---
$\tfrac12 + (\tfrac13{+}\tfrac14) + (\tfrac15{+}\cdots{+}\tfrac18) + \cdots
\ge \tfrac12+\tfrac12+\cdots$ --- is exactly inequality (3.9) written out for
this one series.}
\end{proof}

As a further application of Theorem 3.27, we prove:

\ritem{3.29}{Theorem}
\emph{If $p > 1$,}
\begin{equation}
  \sum_{n=2}^{\infty} \frac{1}{n(\log n)^p} \tag{3.10}
\end{equation}
\emph{converges; if $p \le 1$, the series diverges.}

\emph{Remark:} ``$\log n$'' denotes the logarithm of $n$ to the base $e$
(compare Exercise 1.7); the number $e$ will be defined in a moment (see
Definition 3.30). We let the series start with $n = 2$, since $\log 1 = 0$.

\begin{proof}
The monotonicity of the logarithmic function (which will be discussed in more
detail in Chap.\ 8) implies that $(\log n)_{n=1}^\infty$ increases. Hence
$(1/[n(\log n)^p])_{n=1}^\infty$ decreases, and we can apply Theorem 3.27 to
(3.10); this leads us to the series
\begin{equation}
  \sum_{k=1}^{\infty} 2^k \cdot \frac{1}{2^k(\log 2^k)^p}
  = \sum_{k=1}^{\infty} \frac{1}{(k\log 2)^p}
  = \frac{1}{(\log 2)^p}\sum_{k=1}^{\infty} \frac{1}{k^p}, \tag{3.11}
\end{equation}
and Theorem 3.29 follows from Theorem 3.28.
\end{proof}

This procedure may evidently be continued. For instance,
\begin{equation}
  \sum_{n=3}^{\infty} \frac{1}{n\log n\,\log(\log n)} \tag{3.12}
\end{equation}
diverges, whereas
\begin{equation}
  \sum_{n=3}^{\infty} \frac{1}{n\log n\,(\log(\log n))^2} \tag{3.13}
\end{equation}
converges.

We may now observe that the terms of the series (3.12) differ very little from
those of (3.13). Still, one diverges, the other converges. If we continue the
process which led us from Theorem 3.28 to Theorem 3.29, and then to (3.12) and
(3.13), we get pairs of convergent and divergent series whose terms differ even
less than those of (3.12) and (3.13). One might thus be led to the conjecture
that there is a limiting situation of some sort, a ``boundary'' with all
convergent series on one side, all divergent series on the other side --- at
least as far as series with monotonic coefficients are concerned. This notion
of a ``boundary'' is of course quite vague. The point we wish to make is this:
\emph{No matter how we make this notion precise, the conjecture is false.}
Exercises 3.11(b) and 3.12(b) may serve as illustrations.\kd{A rare editorial
aside from Rudin, and worth taking seriously: there is no slowest divergent
series and no fastest convergent one. Given any convergent $\Sigma a_n$ with
positive terms, Exercise 3.12(b) produces a convergent series with terms
\emph{larger} by an unbounded factor; 3.11(b) does the dual for divergence.
The convergence/divergence frontier is not a curve one can creep up to --- a
genuinely counterintuitive fact, and the reason ``comparison with the worst
convergent series'' can never be a test.}

We do not wish to go any deeper into this aspect of convergence theory, and
refer the reader to Knopp's ``Theory and Application of Infinite Series,''
Chap.\ IX, particularly Sec.\ 41.

\rsection{The Number $e$}

\begin{signpost}[Why $e$ appears in a chapter on series]
Partly because it can --- the series $\sum 1/n!$ is the book's first named sum
--- and partly as a demonstration piece. Theorem 3.31 reconciles the two
classical definitions of $e$ using exactly the $\limsup$/$\liminf$ machinery
just built, and Rudin flags it as ``a good illustration of operations with
limits.'' The error estimate (3.16) then delivers the irrationality of $e$ in
five lines, the chapter's most elegant single proof. Yu Pin's Tsinghua course
takes the definition of $e$ (and eventually $e^{i\pi} = -1$) as the organising
goal of its entire first semester; Rudin gets the number itself this cheaply
because the series technology came first.
\end{signpost}

\ritem{3.30}{Definition}
$\displaystyle e = \sum_{n=0}^{\infty} \frac{1}{n!}$.

Here $n! = 1\cdot2\cdot3\cdots n$ if $n \ge 1$, and $0! = 1$. Since
\begin{align*}
  s_n &= 1 + 1 + \frac{1}{1\cdot2} + \frac{1}{1\cdot2\cdot3} + \cdots +
        \frac{1}{1\cdot2\cdots n}\\
      &< 1 + 1 + \frac12 + \frac{1}{2^2} + \cdots + \frac{1}{2^{n-1}} < 3,
\end{align*}
the series converges, and the definition makes sense.\kn{The comparison is
with the geometric series, of course: $n! \ge 2^{n-1}$ because each factor
past the first is at least $2$. Convergence via 3.24 --- increasing partial
sums, bounded by $3$. So $2 < e < 3$ before anything else is known about it.}
In fact, the series converges very rapidly and allows us to compute $e$ with
great accuracy.

It is of interest to note that $e$ can also be defined by means of another
limit process; the proof provides a good illustration of operations with
limits:

\ritem{3.31}{Theorem}
$\displaystyle\lim_{n\to\infty} \left(1+\frac{1}{n}\right)^n = e$.

\begin{proof}
Let
\[ s_n = \sum_{k=0}^{n} \frac{1}{k!}, \qquad
   t_n = \left(1+\frac{1}{n}\right)^n. \]
By the binomial theorem,
\begin{align*}
  t_n = 1 + 1 &+ \frac{1}{2!}\left(1-\frac1n\right)
        + \frac{1}{3!}\left(1-\frac1n\right)\left(1-\frac2n\right) + \cdots\\
      &+ \frac{1}{n!}\left(1-\frac1n\right)\left(1-\frac2n\right)\cdots
        \left(1-\frac{n-1}{n}\right).
\end{align*}
Hence $t_n \le s_n$, so that
\begin{equation}
  \limsup_{n\to\infty} t_n \le e, \tag{3.14}
\end{equation}
by Theorem 3.19. Next, if $n \ge m$,
\[ t_n \ge 1 + 1 + \frac{1}{2!}\left(1-\frac1n\right) + \cdots +
   \frac{1}{m!}\left(1-\frac1n\right)\cdots\left(1-\frac{m-1}{n}\right). \]
Let $n \to \infty$, keeping $m$ fixed. We get
\[ \liminf_{n\to\infty} t_n \ge 1 + 1 + \frac{1}{2!} + \cdots + \frac{1}{m!}, \]
so that
\[ s_m \le \liminf_{n\to\infty} t_n. \]
Letting $m \to \infty$, we finally get
\begin{equation}
  e \le \liminf_{n\to\infty} t_n. \tag{3.15}
\end{equation}
The theorem follows from (3.14) and (3.15).
\end{proof}

\begin{center}
\begin{tikzpicture}[x=1mm, y=1mm]
  \node[rmbox, text width=30mm] (u0) at (0,24)
    {binomial theorem:\\$t_n \le s_n$ (all $n$)};
  \node[rmbox, text width=28mm] (u1) at (60,24)
    {$\limsup t_n \le e$};
  \draw[rmarrow] (u0) -- node[lbl, anchor=south, inner sep=2.5pt]
    {3.19} (u1);
  \node[rmbox, text width=30mm] (l0) at (0,6)
    {freeze $m$, drop terms:\\$n\to\infty$ by 3.3};
  \node[rmbox, text width=28mm] (l1) at (60,6)
    {$\liminf t_n \ge s_m$,\\then let $m\to\infty$};
  \draw[rmarrow] (l0) -- (l1);
  \node[rmbox, text width=26mm] (m) at (30,-13)
    {$\lim t_n = e$};
  \draw[rmarrow] (u1) -- node[lbl, anchor=west, inner sep=3pt]
    {3.18(c)} (m);
  \draw[rmarrow] (l1) -- (m);
\end{tikzpicture}
\end{center}
\figcap{Two independent one-sided bounds, closed by the squeeze of 3.18(c).
The upper branch may use all $n$ at once; the lower must freeze $m$ first ---
the legal order for exchanging two limits.}

\begin{unpack}[The two-parameter squeeze, dissected]
The proof is a small masterclass in a manoeuvre used constantly from here on.
The obstacle: $t_n$ is a sum of $n$ terms, each \emph{changing with $n$}, so no
single limit theorem applies. The solution: cut the dependence in two.

\smallskip
\textbf{Upper half.} Every bracketed factor is $\le 1$, so term by term
$t_n \le s_n$; take $\limsup$ (which exists unconditionally --- this is why
3.16 was built) and use 3.19.

\smallskip
\textbf{Lower half.} \emph{Freeze} the number of terms at $m$, discard the
rest --- legitimate since all terms are positive --- and let only $n$ run.
Now there are finitely many factors, each $\to 1$, so 3.3 applies and gives
$\liminf t_n \ge s_m$ for each \emph{fixed} $m$. Only then release $m$.

\smallskip
The order of operations is everything: $n$ first with $m$ frozen, $m$ second.
Interchanging two limit processes illegally is the classic analysis error, and
this proof is the book's first demonstration of how to do it legally --- by
one-sided bounds and a squeeze. Example 3.18(c) closes the trap:
$\limsup = \liminf = e$ forces $\lim = e$.
\end{unpack}

The rapidity with which the series $\sum 1/n!$ converges can be estimated as
follows: If $s_n$ has the same meaning as above, we have
\begin{align*}
  e - s_n &= \frac{1}{(n+1)!} + \frac{1}{(n+2)!} + \frac{1}{(n+3)!} + \cdots\\
    &< \frac{1}{(n+1)!}\left\{1 + \frac{1}{n+1} + \frac{1}{(n+1)^2} + \cdots
      \right\} = \frac{1}{n!\,n}
\end{align*}
so that
\begin{equation}
  0 < e - s_n < \frac{1}{n!\,n}. \tag{3.16}
\end{equation}
Thus $s_{10}$, for instance, approximates $e$ with an error less than
$10^{-7}$. The inequality (3.16) is of theoretical interest as well, since it
enables us to prove the irrationality of $e$ very easily.

\ritem{3.32}{Theorem}
\emph{$e$ is irrational.}

\begin{proof}
Suppose $e$ is rational. Then $e = p/q$, where $p$ and $q$ are positive
integers. By (3.16),
\begin{equation}
  0 < q!\,(e-s_q) < \frac{1}{q}. \tag{3.17}
\end{equation}
By our assumption, $q!\,e$ is an integer. Since
\[ q!\,s_q = q!\left(1+1+\frac{1}{2!}+\cdots+\frac{1}{q!}\right) \]
is an integer, we see that $q!\,(e-s_q)$ is an integer.

Since $q \ge 1$, (3.17) implies the existence of an integer between $0$ and
$1$. We have thus reached a contradiction.
\end{proof}

Actually, $e$ is not even an algebraic number. For a simple proof of this, see
page 25 of Niven's book, or page 176 of Herstein's, cited in the Bibliography.

\begin{deeper}[The shape of the irrationality proof]
Compare Chapter 1's proof that $\sqrt2 \notin \Q$: there, rationality forced an
infinite descent; here it forces an integer strictly between $0$ and $1$. Both
are impossibility-of-an-integer arguments, and the second pattern --- scale the
error by exactly the denominator's factorial and watch it land in $(0,1)$ ---
is the template for most irrationality proofs in analysis.

Everything hinges on \emph{how fast} the series converges: (3.16) says the
tail after $s_q$ is smaller than $1/(q!\,q)$, so multiplying by $q!$ leaves
something below $1/q$, still positive. A series converging slower --- like
$\sum 1/n^2$, whose sum $\pi^2/6$ is also irrational --- gives no such
one-line proof. Fast convergence is arithmetic leverage: Exercise 2.3 showed transcendental
numbers \emph{exist} without naming one, and here analysis pins down an
explicit number as irrational.
\end{deeper}
